\documentclass{beamer}
\usetheme{Warsaw}
\usecolortheme{beaver}
\title[Warsaw deck]{A Warsaw Deck with Beaver Colours}
\subtitle{Themes, navigation and blocks}
\author[A. Author]{Ada Author\inst{1} \and Ben Builder\inst{2}}
\institute[Inst.]{\inst{1}First Institute \and \inst{2}Second Institute}
\date{3 October 2026}
\begin{document}
\begin{frame}\titlepage\end{frame}
\begin{frame}{Outline}\tableofcontents\end{frame}
\section{Introduction}
\subsection{Motivation}
\begin{frame}{Motivation}
  \begin{itemize}
    \item Slides should compile exactly as pdflatex compiles them.
    \item Navigation bars show sections and mini frames.
    \item Every theme draws shadows, rounded boxes and gradients.
  \end{itemize}
\end{frame}
\begin{frame}{Blocks}
  \begin{block}{Ordinary block}
    A plain block with a title.
  \end{block}
  \begin{alertblock}{Alert block}
    Something that needs attention: $e^{i\pi}+1=0$.
  \end{alertblock}
  \begin{exampleblock}{Example block}
    An example: $\sum_{k=1}^{n} k = \frac{n(n+1)}{2}$.
  \end{exampleblock}
\end{frame}
\subsection{Layout}
\begin{frame}{Columns}
  \begin{columns}[T]
    \begin{column}{0.48\textwidth}
      \textbf{Left.} A column of text that wraps over several lines so that
      the paragraph builder has something to do.
    \end{column}
    \begin{column}{0.48\textwidth}
      \begin{enumerate}
        \item First
        \item Second
        \item Third
      \end{enumerate}
    \end{column}
  \end{columns}
\end{frame}
\section{Results}
\subsection{Tables}
\begin{frame}{A table}
  \begin{center}
  \begin{tabular}{lrr}
    \hline
    Engine & Pages & Time (ms) \\
    \hline
    pdflatex & 12 & 410 \\
    flashtex & 12 & 4 \\
    \hline
  \end{tabular}
  \end{center}
\end{frame}
\subsection{Math}
\begin{frame}{Equations}
  \begin{align}
    a^2 + b^2 &= c^2 \\
    \int_0^1 x^2\,dx &= \tfrac13
  \end{align}
\end{frame}
\section*{Summary}
\begin{frame}{Summary}
  \begin{itemize}
    \item Warsaw and beaver.
    \item<2-> Overlays in a themed deck.
  \end{itemize}
\end{frame}
\end{document}
